%% =========================================================================
%% Methods (SI) for "Optimal carbon price increased by income distribution
%% effects" -- rewritten to match src/economic_model.py.
%%
%% Substantive changes from draft v1 (postdoc draft):
%%   1.  Added "Income distribution" section.  The draft did not describe
%%       the Lorenz curve at all; the model is built around it.
%%   2.  Clarified what y_net(F) means in the distributional-damage formula:
%%       it is post-climate-damage, post-abatement-tax per-capita income at
%%       population rank F, evaluated one step earlier (one Delta t lag, to
%%       avoid a within-timestep fixed point).
%%   3.  Made the tax base explicit: the equal-pseudo-utility-loss schedule
%%       is applied to damaged income (post-climate, pre-tax), not to gross
%%       income.
%%   4.  Folded the "uniform fractional taxation" subsection into the
%%       general progressive schedule as the tau = 0 (eta_eff = 1) limit;
%%       these are the same closed form in two parameterizations.
%%   5.  Welfare integral: corrected  int u(c(F)) dL(F) -> int u(c(F)) dF.
%%       The dL/dF factor appears only when mapping population rank to per-
%%       capita income, not when aggregating utility across the population.
%%   6.  Made the optimizer's control variable explicit: f(t) = log10(MAC)
%%       in $/tCO2 (and the savings rate s(t)), not the carbon price p_C
%%       directly.  At the optimum MAC = p_C, so the headline result is
%%       unchanged, but the construction matters for reproducibility.
%%   7.  Added one sentence on the one-step lag in omega(F).
%%   8.  Added land emissions and the CO2-equivalent ratio to potential
%%       emissions:  E_pot = sigma * Y_gross * r_{CO2e} + E_land.
%%   9.  Added the Omega_base <= 0.80 cap and stated the budget closure
%%       used for distributional damage when income_dependent_aggregate_
%%       damage = False.
%%
%% Parameter values:
%%   Draft values retained verbatim.  Note for the user: the currently
%%   committed canonical run (json/config_938_100k_6000-3.7_base.json) uses
%%   eta = 1.45 (vs. draft 1.5) and rho = 0.026 (vs. 0.025).  Reconcile
%%   with the postdoc before submission.
%%
%% Code cross-references appear next to each equation, e.g.
%%   % code: economic_model.py:304
%% =========================================================================

\begin{center}
{\Large\bfseries Supporting Information for}\\[6pt]
{\large\bfseries Optimal carbon price increased by income distribution effects}\\[4pt]
{Lamprini Papargyri, Ken Caldeira, and X}
\end{center}

\section*{Model Description}\label{si:derivation}

The model is a single-region, finite-horizon integrated assessment model in the DICE family, extended to resolve an exogenous income distribution and a continuous family of redistributive tax policies.  Each timestep $t \in \{t_0, t_0 + \Delta t, \ldots, t_f\}$ (with $\Delta t = 1$~yr, $t_0 = 2020$, $t_f = 2420$) updates two state variables: capital stock $K(t)$ and cumulative CO$_2$ emissions $E_{\rm cum}(t)$.  Exogenous time functions $A(t), L(t), \sigma(t), \theta_1(t), G(t), r_{\rm CO_2e}(t), E_{\rm land}(t)$ supply total factor productivity, population, carbon intensity of output, the abatement-cost coefficient, the Gini index, the CO$_2$-equivalent-to-CO$_2$ ratio for non-CO$_2$ gases, and land-use emissions.  Two control trajectories are jointly optimized: the logarithmic marginal abatement cost $f(t) \equiv \log_{10}(p_C(t))$ and the savings rate $s(t)$.

The integral over the income distribution that appears in several places below is computed by Gauss--Legendre quadrature with $n_{\rm quad} = 32$ nodes mapped from $[-1,1]$ onto $[0,1]$.

\subsection*{Production and capital}\label{si:economic}

Gross output follows the Cobb--Douglas specification
\begin{equation}\label{eq:si_production}
Y_{\rm gross}(t) \;=\; A(t)\, K(t)^{\alpha}\, L(t)^{1-\alpha},
\qquad \text{equivalently}\qquad
y_{\rm gross}(t) \;=\; A(t)\, \bigl(K(t)/L(t)\bigr)^{\alpha}, % code: economic_model.py:304
\end{equation}
with capital share $\alpha = 0.30$.  Capital accumulates by Euler stepping
\begin{equation}\label{eq:si_capital}
K(t + \Delta t) \;=\; K(t) \;+\; \bigl[\, s(t)\, Y_{\rm net}(t) \;-\; \delta\, K(t) \,\bigr]\, \Delta t, % code: economic_model.py:512,793
\end{equation}
with depreciation rate $\delta = 0.10$~yr$^{-1}$ and $Y_{\rm net}(t)$ defined below in eq.~(\ref{eq:si_ynet_agg}).

\subsection*{Income distribution}\label{si:lorenz}

The income distribution at each $t$ is specified by an exogenous Lorenz curve $L(F; G(t))$ on the population rank $F \in [0,1]$, parameterized by the Gini coefficient $G(t)$.  We use the \emph{empirical} Lorenz form
\begin{equation}\label{eq:si_lorenz_empirical}
L\bigl(F;\, G\bigr) \;=\; \Bigl(1 - \tfrac{G}{G_b}\Bigr)\, F \;+\; \tfrac{G}{G_b}\, L_b(F),
\qquad
L_b(F) \;=\; \sum_{k=0}^{3} w_k\, F^{p_k}, % code: distribution_utilities.py:108-211
\end{equation}
a one-parameter family that linearly interpolates between perfect equality ($G=0$, $L=F$) and an empirical four-term polynomial base curve $L_b$.  The base-curve exponents are $\{p_0, p_1, p_2, p_3\} = \{1.500,\,4.367,\,14.072,\,135.060\}$, the weights are $\{w_0, w_1, w_2, w_3\} = \{0.160,\,0.378,\,0.367,\,0.095\}$ (with $w_0 = 1 - w_1 - w_2 - w_3$), and the base-curve Gini is
\begin{equation}\label{eq:si_Gbase}
G_b \;=\; 1 - 2 \sum_{k=0}^{3} \frac{w_k}{p_k + 1} \;\approx\; 0.681. % code: constants.py:66-71
\end{equation}
Equation~(\ref{eq:si_lorenz_empirical}) is valid for $G \le G_b$.  An alternative single-parameter Pareto--Lorenz form,
\begin{equation}\label{eq:si_lorenz_pareto}
L\bigl(F;\, G\bigr) \;=\; 1 - (1-F)^{1 - 1/a}, \qquad a = \tfrac{1}{2}\bigl(1 + 1/G\bigr), % code: distribution_utilities.py:41-65
\end{equation}
is implemented in the codebase and yields qualitatively similar results.

Per-capita pre-damage income at population rank $F$ is recovered from the slope of the Lorenz curve:
\begin{equation}\label{eq:si_y_of_F}
y_{\rm gross}(F, t) \;=\; y_{\rm gross}(t) \cdot \frac{dL}{dF}\bigl(F;\, G(t)\bigr). % code: economic_model.py:421-434
\end{equation}
Throughout, $F$ refers to the \emph{population rank} (cumulative population mass from below), not the cumulative share of income; the Lorenz slope $dL/dF$ converts between the two.

\subsection*{Climate module}\label{si:climate}

Temperature change is given by the transient climate response to cumulative emissions (TCRE),
\begin{equation}\label{eq:si_tcre}
\Delta T(t) \;=\; k_{\rm climate}\, E_{\rm cum}(t), % code: economic_model.py:290
\end{equation}
with $k_{\rm climate} = 5.687 \times 10^{-13}$~\si{\celsius}/tCO$_2$.  This linear mapping subsumes the carbon cycle: cumulative emissions are taken as the proximate driver of warming, with no atmospheric carbon decay.

Industrial emissions are derived from gross output via an exogenous carbon intensity $\sigma(t)$ and a CO$_2$-equivalent ratio $r_{\rm CO_2e}(t)$, and land emissions are added separately.  Per-capita potential emissions (before abatement) are
\begin{equation}\label{eq:si_epot}
e_{\rm pot}(t) \;=\; \sigma(t)\, y_{\rm gross}(t)\, r_{\rm CO_2e}(t) \;+\; E_{\rm land}(t)/L(t). % code: economic_model.py:351-354
\end{equation}
Realized emissions per capita are $e(t) = (1-\mu(t))\,e_{\rm pot}(t)$, and the cumulative-emission state updates as
\begin{equation}
E_{\rm cum}(t + \Delta t) \;=\; \max\!\bigl\{0,\; E_{\rm cum}(t) + e(t)\,L(t)\,\Delta t\bigr\}. % code: economic_model.py:565,795
\end{equation}

\subsection*{Climate damage}\label{si:damage}

\paragraph{Aggregate damage.}  Aggregate damage is a quadratic function of warming \citep{Barrage2024}:
\begin{equation}\label{eq:si_omega_base}
\Omega_{\rm base}(\Delta T) \;=\; \min\!\Bigl\{0.80,\; \psi_1\, \Delta T + \psi_2\, \Delta T^{2}\Bigr\}, % code: economic_model.py:294, constants.py:76
\end{equation}
with $\psi_1 = 0$ and $\psi_2 = 0.003467$~\si{\celsius}$^{-2}$.  The cap at $0.80$ prevents division by zero and represents the limiting case in which $80\%$ of gross output is lost to climate damage; the model is not intended to be physically informative in that regime.

\paragraph{Distributional damage.}  Climate damage is distributed unequally across income ranks: poorer households experience a larger damage fraction.  The damage fraction at population rank $F$ is parameterized by
\begin{equation}\label{eq:si_omega_F}
\omega(F, t) \;=\; \Omega_{\rm scale}(t)\, y_{\rm net}(F, t)^{-\beta}, % code: economic_model.py:465-470
\end{equation}
where $\beta \in [0,1]$ is the elasticity of the damage fraction with respect to per-capita income, and $y_{\rm net}(F,t)$ is per-capita income at rank $F$ after climate damage and after the abatement tax (see eq.~(\ref{eq:si_post_tax})), evaluated at the previous timestep $t - \Delta t$.  The one-step lag is a numerical convention that avoids a within-timestep fixed point between damage and the post-tax income distribution; the resulting error is $\mathcal{O}(\Delta t)$ and is small for $\Delta t = 1$~yr.

The scale factor $\Omega_{\rm scale}(t)$ is fixed by requiring that the income-weighted average of the per-rank damage fractions equals the temperature-driven aggregate $\Omega_{\rm base}$:
\begin{equation}\label{eq:si_omega_closure}
\Omega_{\rm base}(t) \;=\; \frac{\int_{0}^{1} y_{\rm pretax}(F, t)\, \omega(F, t)\,dF}{\int_{0}^{1} y_{\rm pretax}(F, t)\,dF}, % code: economic_model.py:482-491
\end{equation}
where $y_{\rm pretax}(F, t) = y_{\rm gross}(F, t)\,\bigl(1 - \omega(F, t)\bigr)$ is per-capita income net of damage but before the abatement tax.  Equation~(\ref{eq:si_omega_closure}) implicitly defines $\Omega_{\rm scale}(t)$ at every timestep.

At $\beta = 0$, all households experience the same damage fraction $\Omega_{\rm base}$ (the DICE convention).  At $\beta = 0.5$, a household at half the mean income suffers $\sqrt{2} \approx 1.41$ times the average damage fraction.  At $\beta = 1$, every household loses the same absolute amount, regardless of income.

\paragraph{Aggregate output net of damage.}  Damaged aggregate output is
\begin{equation}\label{eq:si_ydamaged}
Y_{\rm damaged}(t) \;=\; Y_{\rm gross}(t)\,\bigl(1 - \Omega_{\rm base}(t)\bigr). % code: economic_model.py:310
\end{equation}

\subsection*{Abatement}\label{si:abatement}

The optimizer directly controls the trajectory $f(t) = \log_{10}\!\bigl(p_C(t)\bigr)$, the base-10 logarithm of the marginal abatement cost (equivalently, the carbon price at the optimum) measured in \$/tCO$_2$.  Working in log space gives the optimizer roughly uniform sensitivity across the relevant range $p_C \in [10^{-1}, 10^{4}]$~\$/tCO$_2$ (carbon removal through deep abatement).  The control function is represented as a PCHIP interpolant through a small number of control points and is refined iteratively during optimization.

Given $p_C(t)$, the abatement fraction $\mu(t)$ is recovered by inverting the DICE marginal-cost relation
\begin{equation}\label{eq:si_mac}
p_C \;=\; \theta_1\,\mu^{\theta_2-1}
\quad\Longrightarrow\quad
\mu(t) \;=\; \min\!\Bigl\{\mu_{\max}(t),\; \bigl(p_C(t)/\theta_1(t)\bigr)^{1/(\theta_2-1)}\Bigr\}, % code: economic_model.py:357-361
\end{equation}
with cost exponent $\theta_2 = 2.6$, time-varying backstop coefficient $\theta_1(t)$ calibrated to DICE-2023, and an upper bound $\mu_{\max}(t)$ that permits net removal ($\mu > 1$).  Total abatement cost per capita follows the DICE convex form
\begin{equation}\label{eq:si_abatecost}
\mathrm{AbateCost}(t) \;=\; \tfrac{\theta_1(t)}{\theta_2}\, e_{\rm pot}(t)\, \mu(t)^{\theta_2}. % code: economic_model.py:364
\end{equation}
The abatement-cost fraction and aggregate net output are
\begin{equation}\label{eq:si_ynet_agg}
\Lambda(t) \;=\; \frac{\mathrm{AbateCost}(t)}{y_{\rm damaged}(t)},
\qquad
Y_{\rm net}(t) \;=\; \bigl(1 - \Lambda(t)\bigr)\, Y_{\rm damaged}(t). % code: economic_model.py:375-380
\end{equation}

\subsection*{Tax policy: a continuous family of redistributive schedules}\label{si:taxpolicies}

Aggregate abatement cost is collected as a tax on damaged per-capita income.  We parameterize a one-parameter continuous family of progressive schedules by an equity knob $\tau \in [0,1)$, which controls an effective elasticity
\begin{equation}\label{eq:si_eta_eff}
\eta_{\rm eff} \;=\; 1 \;+\; \frac{\tau}{1-\tau}\,(\eta - 1), % code: economic_model.py:400
\end{equation}
where $\eta = 1.5$ is the inequality aversion of the social planner's utility function.

We define a pseudo-utility $U^{\star}(c) = c^{1-\eta_{\rm eff}}/(1-\eta_{\rm eff})$ (with the convention $U^{\star}(c) = \ln c$ at $\eta_{\rm eff} = 1$) and require that each household's pseudo-utility loss due to taxation is a common constant $\mathcal{K}(t)$:
\begin{equation}\label{eq:si_K_def}
U^{\star}\!\bigl(y(F)\bigr) - U^{\star}\!\bigl(c(F)\bigr) \;=\; \mathcal{K}(t)
\qquad \text{for all } F \in [0,1],
\end{equation}
where $y(F) \equiv y_{\rm pretax}(F,t) = y_{\rm gross}(F,t)\bigl(1 - \omega(F,t)\bigr)$ is per-capita damaged but pre-tax income at rank $F$, and $c(F)$ is per-capita post-tax income at rank $F$.  Solving (\ref{eq:si_K_def}) for $c(F)$ yields the closed-form post-tax schedule
\begin{equation}\label{eq:si_post_tax}
c(F, t) \;=\;
\begin{cases}
\Bigl[\, y(F)^{\,1-\eta_{\rm eff}} \;-\; (1-\eta_{\rm eff})\, \mathcal{K}(t) \,\Bigr]^{1/(1-\eta_{\rm eff})}, & \eta_{\rm eff} \ne 1, \\[6pt]
y(F)\,\exp\!\bigl(-\mathcal{K}(t)\bigr), & \eta_{\rm eff} = 1.
\end{cases}
% code: distribution_utilities.py:334-415
\end{equation}
The schedule is monotone in pre-tax income by construction (\,$dc/dy = (c/y)^{\eta_{\rm eff}} > 0$\,), so the ordering of households is preserved by taxation.  The constant $\mathcal{K}(t)$ is fixed by the per-capita budget constraint that the schedule raises exactly the required abatement cost,
\begin{equation}\label{eq:si_budget}
\int_{0}^{1} \Bigl[\, y(F, t) - c(F, t) \,\Bigr]\, dF \;=\; \mathrm{AbateCost}(t), % code: distribution_utilities.py:530-566
\end{equation}
and is computed numerically by a one-dimensional root solve.

\paragraph{Limit cases.}  The progressivity knob spans three regimes:
\begin{itemize}
\item $\tau = 0$ ($\eta_{\rm eff} = 1$): the schedule reduces to the multiplicative form $c(F) = y(F)\,e^{-\mathcal{K}}$.  Every household pays the same fraction $r = 1 - e^{-\mathcal{K}}$ of damaged income.  This is uniform (``flat'') fractional taxation.
\item $\tau = 0.5$ ($\eta_{\rm eff} = \eta$): the pseudo-utility coincides with the planner's CRRA utility, so every household bears the same welfare loss.  Because marginal utility decreases with income ($\eta > 1$), richer households pay larger dollar amounts than poorer ones to deliver the same utility loss.
\item $\tau \to 1$ ($\eta_{\rm eff} \to \infty$): the tax burden concentrates on the highest earners through an increasingly step-like schedule.  The marginal tax rate at rank $F$ is $1 - (c(F)/y(F))^{\eta_{\rm eff}}$, which approaches $1$ for high-income ranks and $0$ for low-income ranks.
\end{itemize}
The single closed form (\ref{eq:si_post_tax}) thus encompasses both the uniform and the highly progressive tax policies as endpoints of a continuous family.

\subsection*{Welfare objective}\label{si:welfare}

Per-capita consumption at rank $F$ is post-tax income net of saving, $c_{\rm cons}(F, t) = (1 - s(t))\, c(F, t)$.  Population-averaged utility at $t$ is
\begin{equation}\label{eq:si_U_t}
U(t) \;=\; \int_{0}^{1} u\!\bigl(c_{\rm cons}(F, t)\bigr)\, dF,
\qquad
u(c) \;=\; \frac{c^{\,1-\eta}}{1-\eta}, % code: economic_model.py:454-472
\end{equation}
with CRRA elasticity $\eta = 1.5$.  The integral is over population mass $dF$, not over the cumulative income share $dL(F)$: the Lorenz slope $dL/dF$ enters only through $y_{\rm gross}(F, t)$ in eq.~(\ref{eq:si_y_of_F}), not through the welfare weighting itself.

The social planner maximizes the population-weighted discounted utility
\begin{equation}\label{eq:si_W}
W \;=\; \sum_{t = t_0}^{t_f} e^{-\rho\,(t - t_0)}\, L(t)\, U(t)\, \Delta t, % code: optimization.py:491-509
\end{equation}
with pure rate of time preference $\rho = 0.025$~yr$^{-1}$, by jointly choosing the control trajectories $\{f(t), s(t)\}$, equivalently $\{p_C(t), s(t)\}$.  In practice each trajectory is represented as a PCHIP interpolant through a small set of control points, and the optimization is solved by the NLopt SBPLX algorithm with progressive refinement of the number of control points.  The savings rate is constrained to $s \in [0.2, 0.4]$ and the log carbon price to $f \in [-1, 4]$.

\subsection*{Welfare cost of abatement and the dependence of $p_C$ on $\tau$}\label{si:welfarecost}

At the social optimum, the first-order condition equates the marginal welfare cost of abatement with the marginal welfare benefit of avoided damage.  The aggregate welfare cost of a marginal dollar of abatement depends on \emph{who pays}: every dollar collected from rank $F$ is weighted by the local marginal utility $u'\!\bigl(c_{\rm cons}(F)\bigr)$.  Under uniform taxation ($\tau = 0$) each household pays the same fraction of damaged income, so the aggregate marginal welfare cost equals the population-average marginal utility.  Under progressive taxation ($\tau > 0$) the burden shifts toward higher-income ranks, where marginal utility is lower, reducing the aggregate marginal welfare cost per dollar of abatement.  In the limit $\tau \to 1$ the marginal welfare cost approaches the marginal utility of the highest earners, which is much smaller than the population average.

Because the optimal carbon price is determined by the marginal-cost--equals--marginal-benefit condition, and because the cost side falls with $\tau$ while the benefit side (avoided damage) is comparatively insensitive to who pays the tax, the optimal carbon price $p_C^{\star}$ \emph{rises} with $\tau$.  This is the income-distribution effect reported in the main text.

\subsection*{Parameter values}\label{si:params}

\begin{center}
\begin{tabular}{l l l}
\hline
\textbf{Symbol} & \textbf{Meaning} & \textbf{Value} \\
\hline
$\alpha$           & Capital share                                  & $0.30$ \\
$\delta$           & Capital depreciation rate                      & $0.10$~yr$^{-1}$ \\
$\psi_1$           & Linear damage coefficient                      & $0$ \\
$\psi_2$           & Quadratic damage coefficient                   & $3.467 \times 10^{-3}$~\si{\celsius}$^{-2}$ \\
$\Omega_{\rm base}^{\max}$ & Cap on aggregate damage fraction       & $0.80$ \\
$k_{\rm climate}$  & TCRE                                           & $5.687 \times 10^{-13}$~\si{\celsius}/tCO$_2$ \\
$\theta_2$         & Abatement-cost exponent                        & $2.6$ \\
$\theta_1(t)$      & Backstop cost                                  & DICE-2023 calibration \\
$\eta$             & Inequality aversion (CRRA)                     & $1.5$ \\
$\rho$             & Pure rate of time preference                   & $0.025$~yr$^{-1}$ \\
$\beta$            & Damage--income elasticity                       & $\in [0,1]$ (varied) \\
$\tau$             & Tax progressivity knob                          & $\in [0,1)$ (varied) \\
$G_b$              & Empirical-Lorenz base Gini                     & $0.681$ \\
$t_0,\,t_f$        & Time horizon                                   & $2020,\,2420$ \\
$\Delta t$         & Timestep                                       & $1$~yr \\
$n_{\rm quad}$     & Gauss--Legendre nodes over $F$                 & $32$ \\
$f_{\min}, f_{\max}$ & Bounds on $\log_{10}(p_C)$                   & $-1,\,4$ \\
$s_{\min}, s_{\max}$ & Bounds on savings rate                       & $0.2,\,0.4$ \\
\hline
\end{tabular}
\end{center}

The time functions $A(t),\, L(t),\, \sigma(t),\, \theta_1(t),\, G(t),\, r_{\rm CO_2e}(t),\, E_{\rm land}(t)$ are taken from a Gompertz / double-exponential fit to the projections of \citet{Barrage2024}, with the Gini trajectory taken from SSP2.

